\documentclass[reqno]{amsart}
\usepackage[T1]{fontenc}
\usepackage{amsmath,amssymb,amsthm}
\usepackage{xr-hyper}
\usepackage[hidelinks]{hyperref}
\externaldocument[rl-]{20261125T010000--rank-lemma-and-zero-estimate__workbook}
\externaldocument[bp-]{20261126T020210--bessel-poincare-series-and-their-gram-operators__workbook}
\externaldocument[m1-]{20261126T051900--unfolding-twice-gram-operators-as-sums-over-rational-points__workbook}
\externaldocument[m2-]{20261126T052000--which-terms-survive-characters-minors-and-the-gain-at-the-level__workbook}
\externaldocument[m3-]{20261126T052200--the-open-cell-siegels-formula-and-poisson-summation__workbook}
\externaldocument[m5-]{20261126T052100--good-forms-and-the-count-of-the-target__workbook}

% Chapter number; \ref{chap} gives it.
\newcounter{chap}
% Chapter numbers, generated from 20261125T012000--grh-workbook__guide.org
% by scripts/make-chapters.py.  Do not edit: change the order of the
% index and rerun the script (./preview-site.sh and CI do this).
%
% Usage in a sheet, after \newcounter{chap}:
%   % Chapter numbers, generated from 20261125T012000--grh-workbook__guide.org
% by scripts/make-chapters.py.  Do not edit: change the order of the
% index and rerun the script (./preview-site.sh and CI do this).
%
% Usage in a sheet, after \newcounter{chap}:
%   % Chapter numbers, generated from 20261125T012000--grh-workbook__guide.org
% by scripts/make-chapters.py.  Do not edit: change the order of the
% index and rerun the script (./preview-site.sh and CI do this).
%
% Usage in a sheet, after \newcounter{chap}:
%   \input{chapters}\chapnumber{<sheet ID>}
\makeatletter
\providecommand\chapters@set[2]{\@namedef{chapters@#1}{#2}}
\chapters@set{20261125T001000}{1}% zeros-to-ramp-sums
\chapters@set{20261125T001100}{2}% inert-primes-and-estimate-b
\chapters@set{20261125T005000}{3}% galois-fourier-and-averaging
\chapters@set{20261125T005500}{4}% columns-from-a-sphere
\chapters@set{20261125T000000}{5}% slope-method
\chapters@set{20261126T030100}{6}% a-twist-for-each-prime
\chapters@set{20261125T003000}{7}% the-chain-and-the-square-root
\chapters@set{20261125T010000}{8}% rank-lemma-and-zero-estimate
\chapters@set{20261125T010500}{9}% poincare-series-and-near-orthogonality
\chapters@set{20261124T223000}{10}% elliptic-curves-for-the-proof
\chapters@set{20261125T009500}{11}% fock-model-and-theta-intertwiners
\chapters@set{20261124T223100}{12}% complex-multiplication-and-inert-reduction
\chapters@set{20261125T008000}{13}% one-cm-determinant
\chapters@set{20261124T214500}{14}% supersingular-theta-divisibility
\chapters@set{20261124T214600}{15}% the-isotropic-plane
\chapters@set{20261125T013700}{16}% pairings-on-a-formal-group-of-height-two
\chapters@set{20261125T013500}{17}% hecke-nondegeneracy-on-the-tree
\chapters@set{20261126T030200}{18}% the-bound-is-attained
\chapters@set{20261125T009000}{19}% kronecker-limit-and-chowla-selberg
\chapters@set{20261125T002000}{20}% definite-theta-lifts
\chapters@set{20261125T015500}{21}% siegel-weil-in-the-smallest-case
\chapters@set{20261125T002100}{22}% definite-theta-lifts-and-norms
\chapters@set{20261125T016500}{23}% hecke-operators-and-theta-lifts
\chapters@set{20261125T016700}{24}% divisibility-along-the-chain
\chapters@set{20261125T004000}{25}% denominators-and-sha
\chapters@set{20261125T004100}{26}% tate-twists-and-sha
\chapters@set{20261126T020430}{27}% the-minus-part-by-gauss-sums
\chapters@set{20261125T014100}{28}% congruence-modules
\chapters@set{20261125T014500}{29}% eisenstein-congruences-and-selmer-classes
\chapters@set{20261125T014700}{30}% pseudocharacters-and-the-length-bound
\chapters@set{20261125T011000}{31}% galois-representations-black-box
\chapters@set{20261126T030300}{32}% the-data
\chapters@set{20261126T020210}{33}% bessel-poincare-series-and-their-gram-operators
\chapters@set{20261126T030400}{34}% the-seesaw-identity
\chapters@set{20261126T030500}{35}% divisibility-of-the-composite-lift
\chapters@set{20261126T030600}{36}% theorem-a-short-forms
\chapters@set{20261126T030700}{37}% the-determinant-of-a-chain
\chapters@set{20261126T030800}{38}% the-cost-of-one-lift
\chapters@set{20261126T030900}{39}% denominators-of-projections-to-lifts
\chapters@set{20261126T031000}{40}% twists-with-few-congruences
\chapters@set{20261126T031100}{41}% theorem-c-and-the-main-inequality
\providecommand\chapnumber[1]{%
  \@ifundefined{chapters@#1}%
    {\PackageError{chapters}{Sheet #1 is not listed in the index}%
      {List it in 20261125T012000--grh-workbook__guide.org\MessageBreak
       and run scripts/make-chapters.py.}}%
    {\edef\thechap{\@nameuse{chapters@#1}}}}
\makeatother
\chapnumber{<sheet ID>}
\makeatletter
\providecommand\chapters@set[2]{\@namedef{chapters@#1}{#2}}
\chapters@set{20261125T001000}{1}% zeros-to-ramp-sums
\chapters@set{20261125T001100}{2}% inert-primes-and-estimate-b
\chapters@set{20261125T005000}{3}% galois-fourier-and-averaging
\chapters@set{20261125T005500}{4}% columns-from-a-sphere
\chapters@set{20261125T000000}{5}% slope-method
\chapters@set{20261126T030100}{6}% a-twist-for-each-prime
\chapters@set{20261125T003000}{7}% the-chain-and-the-square-root
\chapters@set{20261125T010000}{8}% rank-lemma-and-zero-estimate
\chapters@set{20261125T010500}{9}% poincare-series-and-near-orthogonality
\chapters@set{20261124T223000}{10}% elliptic-curves-for-the-proof
\chapters@set{20261125T009500}{11}% fock-model-and-theta-intertwiners
\chapters@set{20261124T223100}{12}% complex-multiplication-and-inert-reduction
\chapters@set{20261125T008000}{13}% one-cm-determinant
\chapters@set{20261124T214500}{14}% supersingular-theta-divisibility
\chapters@set{20261124T214600}{15}% the-isotropic-plane
\chapters@set{20261125T013700}{16}% pairings-on-a-formal-group-of-height-two
\chapters@set{20261125T013500}{17}% hecke-nondegeneracy-on-the-tree
\chapters@set{20261126T030200}{18}% the-bound-is-attained
\chapters@set{20261125T009000}{19}% kronecker-limit-and-chowla-selberg
\chapters@set{20261125T002000}{20}% definite-theta-lifts
\chapters@set{20261125T015500}{21}% siegel-weil-in-the-smallest-case
\chapters@set{20261125T002100}{22}% definite-theta-lifts-and-norms
\chapters@set{20261125T016500}{23}% hecke-operators-and-theta-lifts
\chapters@set{20261125T016700}{24}% divisibility-along-the-chain
\chapters@set{20261125T004000}{25}% denominators-and-sha
\chapters@set{20261125T004100}{26}% tate-twists-and-sha
\chapters@set{20261126T020430}{27}% the-minus-part-by-gauss-sums
\chapters@set{20261125T014100}{28}% congruence-modules
\chapters@set{20261125T014500}{29}% eisenstein-congruences-and-selmer-classes
\chapters@set{20261125T014700}{30}% pseudocharacters-and-the-length-bound
\chapters@set{20261125T011000}{31}% galois-representations-black-box
\chapters@set{20261126T030300}{32}% the-data
\chapters@set{20261126T020210}{33}% bessel-poincare-series-and-their-gram-operators
\chapters@set{20261126T030400}{34}% the-seesaw-identity
\chapters@set{20261126T030500}{35}% divisibility-of-the-composite-lift
\chapters@set{20261126T030600}{36}% theorem-a-short-forms
\chapters@set{20261126T030700}{37}% the-determinant-of-a-chain
\chapters@set{20261126T030800}{38}% the-cost-of-one-lift
\chapters@set{20261126T030900}{39}% denominators-of-projections-to-lifts
\chapters@set{20261126T031000}{40}% twists-with-few-congruences
\chapters@set{20261126T031100}{41}% theorem-c-and-the-main-inequality
\providecommand\chapnumber[1]{%
  \@ifundefined{chapters@#1}%
    {\PackageError{chapters}{Sheet #1 is not listed in the index}%
      {List it in 20261125T012000--grh-workbook__guide.org\MessageBreak
       and run scripts/make-chapters.py.}}%
    {\edef\thechap{\@nameuse{chapters@#1}}}}
\makeatother
\chapnumber{<sheet ID>}
\makeatletter
\providecommand\chapters@set[2]{\@namedef{chapters@#1}{#2}}
\chapters@set{20261125T001000}{1}% zeros-to-ramp-sums
\chapters@set{20261125T001100}{2}% inert-primes-and-estimate-b
\chapters@set{20261125T005000}{3}% galois-fourier-and-averaging
\chapters@set{20261125T005500}{4}% columns-from-a-sphere
\chapters@set{20261125T000000}{5}% slope-method
\chapters@set{20261126T030100}{6}% a-twist-for-each-prime
\chapters@set{20261125T003000}{7}% the-chain-and-the-square-root
\chapters@set{20261125T010000}{8}% rank-lemma-and-zero-estimate
\chapters@set{20261125T010500}{9}% poincare-series-and-near-orthogonality
\chapters@set{20261124T223000}{10}% elliptic-curves-for-the-proof
\chapters@set{20261125T009500}{11}% fock-model-and-theta-intertwiners
\chapters@set{20261124T223100}{12}% complex-multiplication-and-inert-reduction
\chapters@set{20261125T008000}{13}% one-cm-determinant
\chapters@set{20261124T214500}{14}% supersingular-theta-divisibility
\chapters@set{20261124T214600}{15}% the-isotropic-plane
\chapters@set{20261125T013700}{16}% pairings-on-a-formal-group-of-height-two
\chapters@set{20261125T013500}{17}% hecke-nondegeneracy-on-the-tree
\chapters@set{20261126T030200}{18}% the-bound-is-attained
\chapters@set{20261125T009000}{19}% kronecker-limit-and-chowla-selberg
\chapters@set{20261125T002000}{20}% definite-theta-lifts
\chapters@set{20261125T015500}{21}% siegel-weil-in-the-smallest-case
\chapters@set{20261125T002100}{22}% definite-theta-lifts-and-norms
\chapters@set{20261125T016500}{23}% hecke-operators-and-theta-lifts
\chapters@set{20261125T016700}{24}% divisibility-along-the-chain
\chapters@set{20261125T004000}{25}% denominators-and-sha
\chapters@set{20261125T004100}{26}% tate-twists-and-sha
\chapters@set{20261126T020430}{27}% the-minus-part-by-gauss-sums
\chapters@set{20261125T014100}{28}% congruence-modules
\chapters@set{20261125T014500}{29}% eisenstein-congruences-and-selmer-classes
\chapters@set{20261125T014700}{30}% pseudocharacters-and-the-length-bound
\chapters@set{20261125T011000}{31}% galois-representations-black-box
\chapters@set{20261126T030300}{32}% the-data
\chapters@set{20261126T020210}{33}% bessel-poincare-series-and-their-gram-operators
\chapters@set{20261126T030400}{34}% the-seesaw-identity
\chapters@set{20261126T030500}{35}% divisibility-of-the-composite-lift
\chapters@set{20261126T030600}{36}% theorem-a-short-forms
\chapters@set{20261126T030700}{37}% the-determinant-of-a-chain
\chapters@set{20261126T030800}{38}% the-cost-of-one-lift
\chapters@set{20261126T030900}{39}% denominators-of-projections-to-lifts
\chapters@set{20261126T031000}{40}% twists-with-few-congruences
\chapters@set{20261126T031100}{41}% theorem-c-and-the-main-inequality
\providecommand\chapnumber[1]{%
  \@ifundefined{chapters@#1}%
    {\PackageError{chapters}{Sheet #1 is not listed in the index}%
      {List it in 20261125T012000--grh-workbook__guide.org\MessageBreak
       and run scripts/make-chapters.py.}}%
    {\edef\thechap{\@nameuse{chapters@#1}}}}
\makeatother

\ifcsname chapters@20261126T052300\endcsname
  \chapnumber{20261126T052300}
\else
  \renewcommand{\thechap}{?}
\fi
\renewcommand{\thesection}{\thechap.\arabic{section}}

\theoremstyle{definition}
\newtheorem{problem}{Problem}[section]
\theoremstyle{plain}
\newtheorem{theorem}{Theorem}[section]
\theoremstyle{remark}
\newtheorem*{leanremark}{In the Lean formalization}

\newcommand{\C}{\mathbb{C}}
\newcommand{\F}{\mathbb{F}}
\newcommand{\GL}{\mathrm{GL}}
\newcommand{\Gb}{\bar G}
\newcommand{\Rb}{\bar R_0}
\newcommand{\Pco}{{\mathcal P^\circ}}
\newcommand{\Kc}{\mathcal K}
\newcommand{\Uc}{\mathcal U}
\newcommand{\Dc}{\mathcal D}
\newcommand{\Mch}{\mathcal M}
\newcommand{\en}{\mathrm{en}}
\newcommand{\Did}{D_{\mathrm{id}}}
\newcommand{\Dmain}{D_{\mathrm{main}}}
\newcommand{\mt}{\operatorname{mt}}
\newcommand{\rBes}{\rho_{\mathrm{Bes}}}
\newcommand{\tl}[1]{\texttt{#1}}
\setlength{\emergencystretch}{3em}

\title{Positive type and the bootstrap: the second moment}
\author{}
\date{November 26, 2026}

\begin{document}
\maketitle
\refstepcounter{chap}\label{chap}

\paragraph{Prerequisites.}
Hermitian forms and inner products are linear in the first variable; norms of operators are Hilbert--Schmidt norms.  Additive characters of a finite field and their orthogonality \cite[Ch.~5, \S1]{LN}.  From Chapter~\ref{rl-chap}: Problem~\ref{rl-pr:translate-gram}.  From Chapter~\ref{bp-chap}: the Bessel subgroup and the vectors \(v_1,\dots,v_\varrho\), \(z\) (Section~\ref{bp-sec:subgroup}; they belong to the lattice basis of Definition~\ref{bp-def:scaled}), the scaled Poincar\'e map \(\Pco\), the Gram operators \(\Kc_s\) and the covolume \(\nu\) (Definition~\ref{bp-def:scaled}), the group \(\Rb\) and \(\rBes\) (Definition~\ref{bp-def:bessel-rep}), and admissible fillings (Definition~\ref{bp-def:filling}).  From Chapter~\ref{m1-chap}: Theorem~\ref{m1-thm:blocks}, the sizes \(\Did=\#\Rb\cdot D_X\) and \(\Dmain\) (Definition~\ref{m1-def:two-sizes}), Problems~\ref{m1-pr:near-one}, \ref{m1-pr:P-one} and~\ref{m1-pr:positive-type}.  From Chapter~\ref{m2-chap}: frame products and their leading minors \(\langle s\rangle_j\) (Definition~\ref{m2-def:frame}), meeting and the meet norm \(\mt\) (Definition~\ref{m2-def:meet}), open reductions (Definition~\ref{m2-def:open}), fixers (Definition~\ref{m2-def:fixer}), scalar cosets (Definition~\ref{m2-def:scalar-coset}), second moments and link second moments (Definition~\ref{m2-def:moments}), Theorems~\ref{m2-thm:link-bound}, \ref{m2-thm:cut} and~\ref{m2-thm:strip}.  From Chapter~\ref{m3-chap}: Theorem~\ref{m3-thm:open}.

\medskip
Chapter~\ref{bp-chap} deduces Theorem~B from three estimates for the Gram operators \(\Kc_s=\Pco{}^*\rho(s)\Pco\), \(s\in\Gb\).  This sheet proves the second-moment estimate \(\sum_{s\in\Gb}\lVert\Kc_s\rVert^2\le(1+9b^{-4})\Dmain\).  Chapter~\ref{m3-chap} bounds the open reductions that fix given frame vectors; the others, where a leading minor of the frame products vanishes, are many (three sevenths of \(\GL_3(\F_2)\) in the toy below).  How can one show that they carry little of \(s\mapsto\lVert\Kc_s\rVert^2\) without knowing that this function is spread evenly over \(\Gb\)?

\section{Positive type and singular minors}\label{sec:positive}

As in Chapter~\ref{bp-chap}, a real function \(\Theta\) on a finite group \(G\) is \emph{of positive type} if there are vectors \(u_a\), \(a\in G\), of a complex inner product space with \(\langle u_a,u_b\rangle=\Theta(a^{-1}b)\) for all \(a,b\).  A subset of \(G\) \emph{carries} the sum of \(\Theta\) over it.  By Problem~\ref{rl-pr:translate-gram}, if \(\rho\) is a unitary representation of \(G\) and \(L\) a linear map into its space, then \(s\mapsto\lVert L^*\rho(s)L\rVert^2\) is of positive type, with \(u_a=\rho(a)LL^*\rho(a)^*\) in the Hilbert--Schmidt inner product.

\begin{problem}\label{pr:positive}
Let \(\Theta\) be of positive type on \(G\), with vectors \(u_a\).  (a) \(\Theta(s^{-1})=\Theta(s)\) and \(\lvert\Theta(s)\rvert\le\Theta(1)=\lVert u_a\rVert^2\) for all \(s,a\).  (b) For a subgroup \(\Gamma'\subseteq G\), \(\sum_{s\in\Gamma'}\Theta(s)=\lVert\sum_{a\in\Gamma'}u_a\rVert^2/\#\Gamma'\ge0\), and the restriction of \(\Theta\) to \(\Gamma'\) is of positive type.
\end{problem}

Let \(\psi\) be a nontrivial additive character of \(\F_q\); vectors are columns, \(y^t\) is a transpose and \(g^{-t}\) the inverse transpose of \(g\in\GL_n(\F_q)\).

\begin{problem}\label{pr:kernel}
For a \(j\times j\) matrix \(Z\) over \(\F_q\), \(\#\{y\in\F_q^j:y^tZ=0\}=q^{-j}\sum_{y,v\in\F_q^j}\psi(y^tZv)\), and this number is at least \(q\) if \(\det Z=0\).
\end{problem}

\paragraph{The toy.}  In this sheet let \(J\) be the permutation matrix exchanging \(e_1\) and \(e_3\) and fixing \(e_2\), the Gram matrix of the symmetric bilinear form \((x,x')\mapsto x^tJx'\) on \(\F_p^3\).  Take the frame \((e_1,e_2)\): \(e_1\) is isotropic and paired with \(e_3\).  The only proper leading minor of the frame products of \(s\in\GL_3(\F_p)\) is \(e_1^tJse_1=s_{31}\), so the open elements are those with \(s_{31}\ne0\).  An element may fix a frame vector \(x\) in two ways: \(sx=x\), or \(s^{-t}(Jx)=Jx\), that is, \(s\) fixes the linear form \(x'\mapsto x^tJx'\); for \(x=e_1\) the second is \(e_3^ts=e_3^t\).

\begin{problem}\label{pr:gl3-counts}
By machine: in \(\GL_3(\F_p)\) the group, the set \(\{s_{31}=0\}\), the fixer of \(e_1\) and the fixer \(\{s:e_3^ts=e_3^t\}\) of the form \(e_3^t\) have \(168, 72, 24, 24\) elements for \(p=2\) and \(11232, 3456, 432, 432\) for \(p=3\).  Both fixers lie in \(\{s_{31}=0\}\).
\end{problem}

So a weight concentrated on a fixer lives entirely where the minor vanishes.  The next two problems show that this is the only obstruction: a vanishing minor is detected by characters, and the character sums are bounded by the weight of stabilizers.

\begin{problem}\label{pr:character-sum}
Let \(\pi\colon G\to\GL_n(\F_q)\) be a homomorphism, \(\Theta\) of positive type on \(G\), \(y,v\in\F_q^n\) and \(Y\) real.  If the \(s\) with \(\pi(s)v=v\), and the \(s\) with \(\pi(s)^{-t}y=y\), each carry at most \(Y\), then
\[
 \Bigl\lvert\sum_{s\in G}\Theta(s)\,\psi\bigl(y^t\pi(s)v\bigr)\Bigr\rvert\le q^{n/2}Y .
\]
\end{problem}

\begin{problem}\label{pr:det-lemma}
Let \(\pi\) be as in Problem~\ref{pr:character-sum}, \(\Theta\ge0\) of positive type, \(A_0,A_1\) two \(n\times j\) matrices over \(\F_q\) and \(Y\ge0\).  Suppose that for every \(v\in\F_q^j\smallsetminus0\) the \(s\) with \(\pi(s)A_0v=A_0v\), and the \(s\) with \(\pi(s)^{-t}A_1v=A_1v\), each carry at most \(Y\).  Then
\[
 q\sum_{\det(A_1^t\pi(s)A_0)=0}\Theta(s)\le2\sum_{s\in G}\Theta(s)+q^{j}q^{n/2}\,Y .
\]
\end{problem}

\begin{problem}\label{pr:det-toy}
In the toy take \(j=1\), \(A_0=e_1\), \(A_1=Je_1=e_3\), so that \(\det(A_1^tsA_0)=s_{31}\).  (a) For \(\Theta\equiv1\) and \(Y\) the order of a fixer in Problem~\ref{pr:gl3-counts}, Problem~\ref{pr:det-lemma} reads \(144\le471.764\ldots\) for \(p=2\) and \(10368\le29198.213\ldots\) for \(p=3\).  (b) For \(p=2\) the indicator \(\Theta\) of the fixer of \(e_1\) is of positive type; all of its mass \(24\) lies in \(\{s_{31}=0\}\), the least admissible \(Y\) is \(24\), and Problem~\ref{pr:det-lemma} reads \(48\le183.764\ldots\).
\end{problem}

For the uniform weight the elements with \(s_{31}=0\) are a proportion \(3/7\) of \(\GL_3(\F_2)\); for the weight of (b) they are everything.  The function \(s\mapsto\lVert\Kc_s\rVert^2\) is neither, and the proof never replaces its non-open part by a proportion for the uniform measure.  It chooses a subgroup on which stabilizers carry little, and then Problem~\ref{pr:det-lemma} makes the singular part small.  In the toy the lemma gains almost nothing: for \(p=2\) its right side \(471.7\ldots\) exceeds the trivial bound \(2\cdot168\), and for \(p=3\) it is \(29198.2\ldots\) against \(3\cdot11232=33696\).  It gains when \(q^{j+n/2}\,Y\) is small against \(q\sum\Theta\); for the data this is Problem~\ref{pr:exponent} with \(Y=b^{-6k}H\).

\section{The open part and the terminal fixer}\label{sec:fixers}

We work as the proof does: with an admissible filling (Definition~\ref{bp-def:filling}) over a CM field \(E\) with Weil data (Chapter~\ref{m3-chap}); the data provide one, the group \(\Gb=\prod_{w\in\Mch}\GL_r(k_w)\), \(\Mch=\Mch_0\cup\Mch_1\) the chosen places above \(p_0\) and \(p_1\), a unitary representation \(\rho\) of \(\Gb\) on a finite-dimensional complex inner product space, and a linear map \(\Pco\) from the coordinate space with the block identity of Theorem~\ref{m1-thm:blocks}.  In this sheet \(\en(s):=\lVert\Kc_s\rVert^2\), formed with the full Bessel function.  By Problem~\ref{m1-pr:positive-type}, \(\en\) is of positive type, with \(u_a=\rho(a)\Pco\Pco{}^*\rho(a)^*\), and \(\en(s)\le\en(1)\).  The proof also gives \(\en(1)=\lVert\Kc_1\rVert^2\le2D_X\) in this generality (see the remark), by Theorem~\ref{m2-thm:strip} and Problems~\ref{m1-pr:near-one} and~\ref{m1-pr:P-one}, whose hypotheses (\(\Kc_{st}=\Kc_s\rBes(t)\) for \(t\in\Rb\), and \(\rBes\) a unitary representation of \(\Rb\)) it states for every admissible filling with the block identity.  Of the conditions of admissibility we use: \(b\ge1000\); \(b^{10}\le p_0\le b^{11}\) and \(M=p_0p_1\le b^{23}\); \(1000(2\varrho+2)\le k\le f\) and \(\varrho+1\le b\); there are \(f\) chosen places above each \(p_i\), each with residue field of \(p_i\) elements, all distinct, and no conjugate of a chosen place is chosen; \(b^{11kf}\Did\le\Dmain\).  The matrices \(s_w\) have size \(r=k+2\varrho+1\).

We recall Definitions~\ref{m2-def:open}, \ref{m2-def:fixer} and~\ref{m2-def:meet} at a chosen place \(w\) above \(p_0\), with one abbreviation.  In this sheet a \emph{frame vector} at \(w\) is a nonzero vector \(x\) of the span of the frame modulo \(w\), with a side: \(s\) fixes it on the side of \(w\) if \(s_wx=x\), and on the side of \(\bar w\) if \(s_w^{-t}\) fixes \(Jx\), where \(J\) is the Gram matrix of the lattice basis modulo \(w\).  The fixer of frame vectors given at a set \(\Dc\subseteq\Mch_0\) of such places is the subgroup of the \(s\) fixing each.  An element \(s\) is \emph{open} if at some \(w\in\Mch_0\) the leading minors \(\langle s\rangle_1,\dots,\langle s\rangle_\varrho\) of its frame products are nonzero.  With \(A_0\) the matrix whose columns are the first \(j\) vectors \(v_1,\dots,v_j\) of the frame and \(A_1=JA_0\), the minor \(\langle s\rangle_j\) at \(w\) is \(\det(A_1^ts_wA_0)\).  An element \emph{meets} at \(w\) (or at \(\bar w\)) if modulo that place the span of the frame and that of its image share a nonzero vector, and \(\mt(s)\) is the product of \(\#k_w\) over the places \(w\) and \(\bar w\), \(w\in\Mch\), at which \(s\) meets; an element that fixes a frame vector on a side meets at that place.  An element of \(\Rb\) fixes \(v_1\) modulo every chosen place: \(\Uc\) and \(U(X_0)\) fix \(v_1\) (Section~\ref{bp-sec:subgroup}), and \(v_1\) is a vector of the lattice basis.  For \(s\notin\Rb\) we have \(\rBes(s)=0\), and \(\en(s)\) is the second moment at \(s\) of the terms outside the Bessel subgroup.

Three statements are used, for sets \(\mathcal B\subseteq\Gb\).  Theorem~\ref{m2-thm:link-bound}: if \(z_0\) is a positive integer with \(z_0\le\mt(s)\) for all \(s\in\mathcal B\), the link second moment of the terms with \(\operatorname{disc}\ne0\) over \(\mathcal B\) is at most
\[
 (b^{2f}/z_0)^k\,\Dmain\,\nu^2 .
\]
Theorem~\ref{m2-thm:cut}: if \(\mathcal B\) avoids the scalar cosets and that link second moment is at most \(C_0\nu^2\), \(C_0\ge0\), then the second moment over \(\mathcal B\) of the terms outside the Bessel subgroup, for the full Bessel function, has square root at most
\[
 \sqrt{C_0}+b^{-15}\sqrt{\Did}.
\]
Theorem~\ref{m3-thm:open}: for a set \(\Dc\) of \(h\) chosen places above \(p_0\) with a frame vector at each, the link second moment of the terms with \(\operatorname{disc}\ne0\) over the open elements of their fixer is at most
\[
 (1+b^{-4})\,b^{-6kh}\,\Dmain\,\nu^2 .
\]
In this sheet \(y_h:=b^{-6kh}\Dmain\) for \(0\le h\le f\); since \(b^{11kf}\Did\le\Dmain\), we have \(b^{5kf}\Did\le y_h\).

\begin{problem}\label{pr:exponent}
\(p_0^j\sqrt{p_0^{\,r}}\le b^{6k}\) for \(0\le j\le\varrho\).
\end{problem}

\begin{problem}\label{pr:non-open}
Let \(\Gamma'\subseteq\Gb\) be a subgroup, \(H:=\sum_{s\in\Gamma'}\en(s)\) (in this sheet), and \(w\) a chosen place above \(p_0\).  Suppose that for every frame vector at \(w\), on either side, the elements of \(\Gamma'\) fixing it carry at most \(b^{-6k}H\).  Then the elements of \(\Gamma'\) that are not open carry at most \(b^{-4}H\).
\end{problem}

\begin{problem}\label{pr:scalar-not-open}
If \(s_wv_1\in k_wv_1\) at a chosen place \(w\) above \(p_0\), then \(\langle s\rangle_1=0\) at \(w\).  Hence no \(s\in\Gb\) with \(s_w=\zeta_wt_w\) for every \(w\in\Mch\), for some \(t\in\Rb\) and some integral finite adele \(\zeta\) with residues \(\zeta_w\), is open; in particular the open elements avoid the scalar cosets and lie outside \(\Rb\).
\end{problem}

\begin{problem}\label{pr:open-part}
Let \(\Gamma'\) be the fixer of frame vectors at a set of \(h\) chosen places above \(p_0\).  The open elements of \(\Gamma'\) carry at most \((1+4b^{-4})\,y_h\).
\end{problem}

The terminal case fixes vectors at all \(f\) places above \(p_0\).

\begin{problem}\label{pr:terminal-sizes}
(a) \((b^{2f}/p_0^f)^k\,\Dmain\le y_f/8\).  (b) \(8M^f\Did\le y_f\).  (c) At most \(M^f\#\Rb\) elements of \(\Gb\) have the form \(s_w=\zeta_wt_w\) of Problem~\ref{pr:scalar-not-open}.
\end{problem}

\begin{problem}\label{pr:terminal}
Let \(\Gamma'\) be the fixer of frame vectors at all \(f\) chosen places above \(p_0\).  Then \(\mt(s)\ge p_0^f\) for \(s\in\Gamma'\), and \(\Gamma'\) carries at most \(y_f\).
\end{problem}

\section{The bootstrap}\label{sec:bootstrap}

Problem~\ref{pr:non-open} needs a subgroup on which stabilizers carry little, and fixing a vector costs a factor \(b^{-6k}\) in Problem~\ref{pr:open-part}.  For a set \(\Dc\) of chosen places above \(p_0\) and frame vectors at its places, let \(\Gamma_{\Dc}\) be their fixer (the vectors understood); \(\Gamma_\varnothing=\Gb\).  The proof maximizes the \emph{objective} \(b^{6k\#\Dc}\sum_{s\in\Gamma_{\Dc}}\en(s)\) over the finitely many choices, the empty one included.

\begin{problem}\label{pr:bootstrap-toy}
In the toy with \(p=2\), \(\Theta\equiv1\), one place, and the factor \(b^{6k}\) replaced by \(c\): every frame vector, on either side, has a fixer of \(24\) elements, and the objective at (no place, the place) is \((168,24c)\); it is \((168,96)\) for \(c=4\) and \((168,192)\) for \(c=8\).  Either choice can be the maximum.
\end{problem}

\begin{problem}\label{pr:bootstrap}
Let \(\Dc\) with its frame vectors maximize the objective, \(h:=\#\Dc\) and \(H:=\sum_{s\in\Gamma_{\Dc}}\en(s)\).  If \(h<f\) and \(w\notin\Dc\) is a chosen place above \(p_0\), then the elements of \(\Gamma_{\Dc}\) fixing a frame vector at \(w\), on either side, carry at most \(b^{-6k}H\), and
\[
 H\le(1+4b^{-4})\,y_h+b^{-4}H .
\]
\end{problem}

\begin{problem}\label{pr:absorb}
If \(H\le(1+4\delta)Y+\delta H\) with \(0\le\delta\le1/10\) and \(Y\ge0\), then \(H\le(1+9\delta)Y\).  For \(\delta=1/16\): \((1+4\delta)/(1-\delta)=4/3\le25/16=1+9\delta\).
\end{problem}

\begin{problem}\label{pr:whole}
With \(\Dc\) as in Problem~\ref{pr:bootstrap}, \(H\le(1+9b^{-4})\,y_h\); hence \(\sum_{s\in\Gb}\en(s)\le(1+9b^{-4})\,\Dmain\).
\end{problem}

Problems~\ref{pr:positive}, \ref{pr:kernel}, \ref{pr:character-sum}, \ref{pr:det-lemma} and \ref{pr:exponent}--\ref{pr:whole}, with Theorems~\ref{m2-thm:link-bound}, \ref{m2-thm:cut} and~\ref{m3-thm:open} and the two properties of \(\en\) recalled in Section~\ref{sec:fixers} (positive type, and \(\en(1)\le2D_X\)), prove the following.

\begin{theorem}[the second moment]\label{thm:second-moment}
Let an admissible filling over a CM field with Weil data be given, \(\rho\) a unitary representation of \(\Gb\) on a finite-dimensional complex inner product space, and \(\Pco\) a linear map from the coordinate space satisfying the block identity of Theorem~\ref{m1-thm:blocks}.  Then
\[
 \sum_{s\in\Gb}\lVert\Kc_s\rVert^2\le(1+9b^{-4})\,\Dmain .
\]
In particular this holds for the data, with \(\Pco\) the scaled Poincar\'e map and \(\rho\) the representation of Chapter~\ref{bp-chap}.
\end{theorem}

What has been proved is the second moment over the whole group, from the link bound and the cut, which Chapter~\ref{m2-chap} reads from the proof, and the open-cell bound, whose steps Chapter~\ref{m3-chap} models with statements read from the proof.  With the count of Theorem~\ref{m5-thm:count} it gives the hypothesis of Theorem~\ref{rl-thm:spread} on the second moment (Problem~\ref{m5-pr:spread-hyp}); the other hypothesis, the near isometry of the projected map, is Theorem~\ref{m5-thm:good}.

\begin{leanremark}
\raggedright In the Lean proof \cite{GL} the theorem is \tl{Filling.HasAmplitudeBlocks.sum\_hilbertSchmidtNorm\_sq\_le}, and for the data \tl{Data.sum\_hilbertSchmidtNorm\_scaledPoincare\_sq\_le}; it is deduced from \tl{Filling.sum\_secondMoment\_le}, which bounds the sum over \(\Gb\) of the second moments at \(s\) of all terms, for an admissible filling with Weil data, assuming that this function of \(s\) is of positive type and at most \(2D_X\) at \(s=1\), by the maximization of Section~\ref{sec:bootstrap}.  The bound \(\en(1)\le2D_X\) under the block identity, for any representation, is \tl{Filling.HasAmplitudeBlocks.hilbertSchmidtNorm\_adjoint\_comp\_sq\_le}.  Problem~\ref{pr:det-lemma} is \tl{card\_mul\_sum\_det\_eq\_zero\_le}, whose proof bounds only the real part of the sum of Problem~\ref{pr:character-sum}; Problem~\ref{pr:non-open} is the case \(\en\) of \tl{Filling.sum\_not\_isOpenElement\_le}, stated there for every function \(\ge0\) of positive type; the second sentence of Problem~\ref{pr:scalar-not-open}, up to ``in particular'', is \tl{Filling.InScalarCoset.not\_isOpenElement}; Problem~\ref{pr:open-part} is \tl{Filling.secondMoment\_fixer\_isOpenElement\_le}; the bound for the fixer in Problem~\ref{pr:terminal} is \tl{Filling.secondMoment\_fixer\_univ\_le}, which assumes \(\en\le2D_X\); the first sentence of Problem~\ref{pr:absorb} is \tl{Filling.le\_one\_add\_nine\_mul\_of\_le\_add}.
\end{leanremark}

\section*{Hints}

\subsection*{\S\ref{sec:positive}}
\begin{description}
\item[\ref{pr:positive}] \(\langle u_a,u_a\rangle=\Theta(1)\) and Cauchy--Schwarz.  Expand \(\lVert\sum_au_a\rVert^2\) and group the pairs \((a,b)\) by \(a^{-1}b\).
\item[\ref{pr:kernel}] \(\sum_{v}\psi(z^tv)\) is \(q^j\) if \(z=0\) and \(0\) otherwise; a nonzero \(y\) in the kernel gives a line.
\item[\ref{pr:gl3-counts}] Run through the \(p^9\) matrices.
\item[\ref{pr:character-sum}] Write the sum as an inner product of two vectors of \(\C^{\F_q^n}\otimes\operatorname{span}(u_a)\), one built from point masses, the other from the normalized characters \(\xi\mapsto q^{-n/2}\psi(z^t\xi)\); then Cauchy--Schwarz.
\item[\ref{pr:det-lemma}] Bound \(q\) times the indicator of \(\det Z_s=0\), \(Z_s:=A_1^t\pi(s)A_0\), by Problem~\ref{pr:kernel}; sum against \(\Theta\).  Separate the pairs with \(y=0\) or \(v=0\), and use \(y^tZ_sv=(A_1y)^t\pi(s)(A_0v)\) for the others.
\item[\ref{pr:det-toy}] For (b), \(u_a:=\delta_{ae_1}\).
\end{description}

\subsection*{\S\ref{sec:fixers}}
\begin{description}
\item[\ref{pr:exponent}] Show \(11(2j+r)\le12k\), using \(p_0\le b^{11}\), \(r=k+2\varrho+1\), \(1000(2\varrho+2)\le k\).
\item[\ref{pr:non-open}] A non-open element has \(\langle s\rangle_j=0\) at \(w\) for some \(1\le j\le\varrho\).  For each \(j\), Problem~\ref{pr:det-lemma} for \(\Gamma'\), \(\pi(s)=s_w\), \(\Theta=\en\) and \(Y=b^{-6k}H\): \(A_0v\) and \(A_1v=J(A_0v)\) are a frame vector and its form.  Then Problem~\ref{pr:exponent} and \(\varrho+1\le b\).
\item[\ref{pr:scalar-not-open}] \(v_1\) is isotropic.
\item[\ref{pr:open-part}] By Problem~\ref{pr:scalar-not-open}, Theorems~\ref{m3-thm:open} and~\ref{m2-thm:cut} apply to the open elements of \(\Gamma'\).  Then \(\sqrt{1+\epsilon}\le1+\epsilon/2\), \(b^{-15}\le b^{-4}/2\) and \(\Did\le y_h\).
\item[\ref{pr:terminal-sizes}] (a) \(p_0\ge b^{10}\).  (b) \(23f+1\le5kf\).  (c) Only the residues of \(\zeta\) at the \(2f\) chosen places matter.
\item[\ref{pr:terminal}] The pairs (place, side) at which the fixed vectors sit are distinct places above \(M\).  Split \(\Gamma'\) into the elements of Problem~\ref{pr:terminal-sizes}(c) and the rest; use \(\en\le2D_X\), Theorems~\ref{m2-thm:link-bound} with \(z_0=p_0^f\) and~\ref{m2-thm:cut}, and \((\sqrt C+\sqrt{C'})^2\le2C+2C'\).
\end{description}

\subsection*{\S\ref{sec:bootstrap}}
\begin{description}
\item[\ref{pr:bootstrap-toy}] A nonzero vector of \(\F_2^3\), or a nonzero form, has \(7\) images.
\item[\ref{pr:bootstrap}] Adding \(w\) with one frame vector replaces \(\Gamma_{\Dc}\) by its intersection with the fixer of that vector; compare the objectives.  Then Problems~\ref{pr:positive}(b), \ref{pr:non-open} and~\ref{pr:open-part}.
\item[\ref{pr:absorb}] \((1-\delta)(1+9\delta)-(1+4\delta)=\delta(4-9\delta)\).
\item[\ref{pr:whole}] If \(h=f\), Problem~\ref{pr:terminal}; otherwise \(\delta=b^{-4}\).  Then compare with \(\Dc=\varnothing\).
\end{description}

\begin{thebibliography}{LN}
\bibitem[GL]{GL} L.~Sernau, \emph{GRH for Dirichlet \(L\)-functions, in Lean~4}, a shorter rewriting of the Lean proof, text by Claude, \url{https://github.com/lsernau/grh-lean}, commit \texttt{72fead1}, readable edition \texttt{grh/}.
\bibitem[LN]{LN} R.~Lidl and H.~Niederreiter, \emph{Finite Fields}, 2nd ed., Encyclopedia of Mathematics and its Applications \textbf{20}, Cambridge University Press, 1997.
\end{thebibliography}
\end{document}
